This chapter presents the thesis's central design contribution: a four-stage pipeline that turns a
just-answered query, a user's own query history, and their behavioral profile
(Chapter~\ref{ch:behavior}) into three follow-up query suggestions.

\section{Overview}
\label{sec:rec-overview}

The central research question --- can conditioning on real behavior, not just query history,
produce measurably different recommendations --- is operationalized directly in this chapter's
design: the candidate-generation prompt (Section~\ref{sec:rec-generation}) explicitly presents the
LLM with two separate signals, framed as potentially \emph{different} things, not two views of the
same preference; and the ranking stage (Section~\ref{sec:rec-ranking}) makes this signal
load-bearing rather than decorative, rewarding candidates that align with the behavioral profile's
real top categories, not merely with the user's own past phrasing. All four stages are composed
inside the single extension point Section~\ref{sec:arch-api} established,
\code{handle\_recommend()}, requiring no change to the API contract or the development harness.

\begin{figure}[H]
\centering
\begin{adjustbox}{max width=\textwidth, max totalheight=0.88\textheight, center}
\begin{tikzpicture}[node distance=1.0cm and 2.2cm]

    \node[block, fill=colSource] (qlog) {Query log\\[2pt]\footnotesize \code{query\_log.db}};
    \node[block, fill=colSource, right=of qlog] (bprof) {Behavioral profile\\[2pt]\footnotesize Ch.~\ref{ch:behavior}};

    \node[block, fill=colGreen, below=of qlog] (retrieve) {4a --- retrieval\\[2pt]\footnotesize top-$k$ similar past queries, embedding cosine};

    \node[block, fill=colOrange, below=of retrieve] (generate) {4b --- generation\\[2pt]\footnotesize LLM prompt (query + history + divergence) $\to$ pool of 5};

    \node[block, fill=colPurple, below=of generate] (validate) {4c --- validation\\[2pt]\footnotesize schema-grounding cosine-similarity floor};

    \node[block, fill=colPink, below=of validate] (rank) {4d --- ranking\\[2pt]\footnotesize history-similarity vs.\ behavioral-alignment, MMR diversify};

    \node[smallblock, fill=colPink, below=of rank] (out) {3 suggestions};

    \draw[arrow] (qlog) -- (retrieve);
    \draw[arrow] (retrieve) -- (generate);
    \draw[arrow] (bprof) -- (generate);
    \draw[arrow] (generate) -- (validate);
    \draw[arrow] (validate) -- (rank);
    \draw[arrow] (retrieve.east) .. controls +(2.5,0) and +(2.5,0.8) .. (rank.east);
    \draw[arrow] (bprof.east) .. controls +(2.2,0) and +(2.2,-6.5) .. (rank.east);
    \draw[arrow] (rank) -- (out);

\end{tikzpicture}
\end{adjustbox}
\caption{The four-stage recommender core. The behavioral profile (right) feeds both the generation
prompt and the ranking score directly --- the two concrete points where the thesis's
stated-vs-actual signal shapes the final output, not just the initial candidate pool.}
\label{fig:rec-pipeline}
\end{figure}

\section{Retrieval}
\label{sec:rec-retrieval}

The retrieval stage implements the retrieval half of RA-GQR~\cite{bacciu2024}: a user's query
history is pulled from the query log (Section~\ref{sec:arch-api}), exact duplicates of the current
query are excluded, the current query and every remaining history entry are embedded with the same
sentence-embedding model already used elsewhere in the NL2SPARQL pipeline for class/property
relevance extraction, and candidates are ranked by cosine similarity. The top-$k$ retrieved queries
become the few-shot context for Section~\ref{sec:rec-generation}'s prompt, in place of generically
hand-curated examples. If a user has no other logged queries, the function returns an empty list;
the generation stage still functions correctly, consistent with the original GQR premise that the
technique does not require historical data to work at all.

This stage was verified directly: seeding one user's query log with four varied movie queries, and
querying with \emph{``Show me exciting action films,''} correctly ranked \emph{``Action movies''}
(0.79 cosine) and \emph{``Best action movies with high ratings''} (0.765) above \emph{``Romantic
comedies''} (0.371), excluding unrelated seeded queries entirely.

\section{Candidate Generation}
\label{sec:rec-generation}

A prompt is assembled from three inputs: the current query; the retrieved similar past queries,
rendered as a bulleted list; and a compact rendering of the behavioral divergence
(Section~\ref{sec:behavior-divergence}), explicitly framed as \emph{``what this user says they
want''} versus \emph{``what this user actually engages with, which may DIFFER.''} The LLM is asked
for a pool of \textbf{five} candidates, deliberately more than the three ultimately returned, so
that the validation and ranking stages have real material to filter and diversify from.

Generation reuses the same multi-provider LLM abstraction already used elsewhere in the codebase,
so the recommender can run against any configured provider --- cloud or local --- without code
changes, only an environment variable and that provider's credentials. In practice, both configured
cloud providers (Gemini and OpenAI) turned out to be unusable during this thesis for account and
billing reasons unrelated to the code: one had no billing configured, the other had a zero free-tier
quota on its project. Rather than block on acquiring new credentials, a local, containerized Ollama
model (\code{llama3.1:8b}) was integrated as an additional provider, requiring no API key and no
quota, and was used for all generation reported in this thesis, including the evaluation in
Chapter~\ref{ch:eval}. This is explicitly a validation of the \emph{pipeline}, not of the
\emph{provider} originally envisioned; the abstraction was already provider-agnostic before this
substitution, so a funded cloud key could be swapped in later with no code change. At this
substitution's local, CPU-only configuration, one generation call takes roughly 18--22 seconds ---
acceptable for offline evaluation, but worth flagging as a real deployment-latency concern distinct
from suggestion quality.

\section{Schema-Grounding Validation}
\label{sec:rec-validation}

\subsection{Motivation and original design}

Section~\ref{sec:rel-kgllm} establishes the motivation directly: LLM-generated content referencing
structured knowledge is a documented hallucination risk, and the baseline system's own measured
SPARQL-correctness rate (70\%, Section~\ref{sec:rel-nl2sparql}) is a concrete reminder that this
risk is not merely theoretical in this exact class of system. The original design intent was to run
each candidate query's implied entities and classes through the NL2SPARQL pipeline's own
class-relevance and property-relevance evaluators, confirming a candidate is answerable against the
currently loaded domain schema before it is ever shown to a user.

\subsection{Empirical calibration and a negative result}

Two implementations of this idea were built and empirically rejected before a third, deliberately
weaker version shipped.

The first, reusing the pipeline's own class-relevance evaluator directly, uses relative
thresholding --- keeping every class scoring at or above one standard deviation below the mean
across the schema's own classes. With only 7--11 classes per domain, this reliably selects at least
one class for \emph{any} input, including completely unrelated text: evaluating \emph{``Best pizza
recipe''} against the Tourism schema returned a non-empty class list.

The second computed raw cosine similarity between each candidate and every class's embedded
description, requiring the maximum score across all classes to clear a fixed floor. This was
calibrated against control (clearly off-topic) and probe (genuinely in-domain) phrases
(Table~\ref{tab:rec-calibration}).

\begin{table}[H]
\centering
\begin{tabular}{lclr l}
\toprule
\textbf{Query} & \textbf{Domain} & \textbf{Max cosine sim.} & \textbf{Best-matching class} \\
\midrule
Top rated action movies from the 2010s        & movie    & 0.275 & Movie \\
Movies directed by Christopher Nolan          & movie    & 0.302 & Movie \\
Documentaries produced in Canada              & movie    & 0.288 & Country \\
Best pizza recipe \emph{(control)}            & movie    & 0.175 & SubscriptionPlan \\
How to fix a flat tire \emph{(control)}       & movie    & 0.094 & Country \\
\textbf{Churches in Verona}                   & tourism  & 0.163 & Location \\
\textbf{Roman Churches in Verona}             & tourism  & 0.154 & Location \\
\textbf{Free Churches in Verona}              & tourism  & \textbf{0.092} & Location \\
Events happening this weekend                 & tourism  & 0.401 & Event \\
Art categories available                      & tourism  & 0.601 & ArtCategory \\
Best pizza recipe \emph{(control)}            & tourism  & 0.064 & Tour \\
\textbf{How to fix a flat tire (control)}     & tourism  & \textbf{0.112} & Tour \\
Latest news about the stock market \emph{(control)} & tourism & 0.004 & --- \\
\bottomrule
\end{tabular}
\caption{Absolute cosine-floor calibration data. The Movie domain separates cleanly, but the
Tourism domain does not: the supervisor's own canonical example, \emph{``Free Churches in
Verona,''} scores below an off-topic control phrase.}
\label{tab:rec-calibration}
\end{table}

The Movie domain separates cleanly (in-domain floor $\approx 0.275$, control ceiling $\approx
0.175$). \textbf{The Tourism domain does not}: \emph{``Free Churches in Verona''} --- one of the
supervisor's own three canonical example suggestions, quoted verbatim in the original scoping email
(Section~\ref{sec:intro-motivation}) --- scores $0.092$, below the off-topic control phrase
\emph{``How to fix a flat tire''} at $0.112$. No single threshold can accept the former while
rejecting the latter. A $z$-score variant, relative to each query's own mean and standard deviation
across classes, was also tried and rejected for the same underlying reason: with only 7--11 classes,
the maximum of such a small sample is reliably 1.3--2.3 standard deviations above the mean
regardless of whether the query is on-topic, confirmed empirically --- \emph{``Best pizza recipe''}
against the Movie schema scored a \emph{higher} $z$-score than several genuinely in-domain movie
queries.

The root cause is structural, not a tuning failure: both domains' machine-readable schema class
descriptions are generated purely from SHACL cardinality constraints (e.g.\ \emph{``Art has exactly
1 Art Class ID''}), with no natural-language descriptions and no example instance values in either
source ontology. The Tourism ontology's classes carry essentially no text a general-purpose sentence
embedding could relate to English words like ``church'' --- that correspondence only exists via the
literal Italian category value \emph{``Chiese''} in the live database, which is not present anywhere
in the schema metadata being embedded.

\subsection{What shipped instead}

Given this finding, the shipped validation step uses the absolute-cosine-floor mechanism with a
threshold of $0.05$ --- low enough to reject only genuinely degenerate candidates (near-zero or
negative similarity to every class in the schema, such as \emph{``Latest news about the stock
market''} at $0.004$) while admitting everything else, including borderline-but-valid phrasings like
\emph{``Free Churches in Verona.''} This is explicitly documented as a coarse sanity filter, not the
precision hallucination guard originally intended, and is re-verified against four cases:
\emph{``Free Churches in Verona''} (tourism) and \emph{``Roman Churches in Verona''} (tourism) both
correctly pass; \emph{``Top rated action movies''} (movie) correctly passes; \emph{``Latest news
about the stock market''} (tourism) correctly fails.

This negative result is treated here as a genuine, reportable contribution
(Section~\ref{sec:intro-contrib}) rather than a silently shipped limitation: automatic
schema-grounding of natural-language candidate suggestions, at the class-description level, was not
achievable with acceptable precision given how sparse this project's source ontologies are. A
stronger version would require enriching class descriptions with real catalog/instance values pulled
from the live VKG endpoints --- a concrete, scoped direction for future work
(Section~\ref{sec:conclusion-future}).

\section{Ranking and Diversification}
\label{sec:rec-ranking}

The final stage scores surviving candidates along two axes echoing QRMOCCGA's two-objective framing
(Section~\ref{sec:rel-qrmoccga}), each normalized to $[0,1]$: history-similarity, the maximum
Jaccard word overlap between a candidate and any of the user's retrieved similar past queries; and
behavioral-alignment, the fraction of the behavioral profile's real top interest categories
(Chapter~\ref{ch:behavior}) whose words appear in the candidate text. This second axis is the
concrete place the thesis's central signal enters ranking: a candidate that textually references
what the user \emph{actually watches}, not just what they typed, scores higher here. A weighted sum
of the two axes ranks the candidate pool; a greedy, Maximal-Marginal-Relevance-style loop then
repeatedly selects the highest-scoring remaining candidate, penalized by its word overlap with
candidates already selected, so the final three suggestions are not near-duplicates of one another.

This stage was verified with a hand-built pool of five candidates (one near-duplicate of query
history, two behaviorally-aligned, one off-axis, one more history-similar): it correctly picked the
history-similar candidate first, then diversified into the two behaviorally-aligned candidates
rather than a near-duplicate of the first pick. Reimplementing QRMOCCGA's cooperative
co-evolutionary genetic algorithm itself was explicitly out of scope --- only its feature ideas were
borrowed, as planned in Section~\ref{sec:rel-qrmoccga}.

\section{End-to-End Verification}
\label{sec:rec-verification}

The full chain was verified at three levels of fidelity. First, with a mocked LLM returning fixed
candidates, the real production entry point --- not a reimplementation --- was called for the
user profiled in Section~\ref{sec:behavior-results} with the query \emph{``Show me exciting action
films.''} The result:

\begin{lstlisting}[language=, caption={Mocked-LLM end-to-end output, ranked top to bottom.}]
War documentaries produced in the USA
Crime thrillers with high ratings
Adventure movies set during wartime
\end{lstlisting}

The top-ranked suggestion is the one most strongly aligned with this user's \emph{actual} top
interest categories (\code{Movie, USA, Canada, Adventure, War}) rather than a plain paraphrase of
the stated query about ``action films.''

Second, over real HTTP with a placeholder API key, the server log confirmed the full path executed
--- log, retrieval, a generation attempt that failed cleanly on the invalid key, caught and logged
as a warning rather than a crash --- and the response was \code{200 OK} with an empty suggestion
list, confirming the pipeline degrades gracefully rather than failing the request when generation is
unavailable.

Third, and most significantly, against genuine Ollama-backed model output (Section~\ref{sec:rec-generation}),
the same request produced:

\begin{lstlisting}[language=, caption={End-to-end suggestions from genuine \code{llama3.1:8b}
output, via the real production code path.}]
{
  "suggestions": [
    "War movies with high ratings on IMDB",
    "Movies that combine action and adventure genres",
    "Films starring actors from Canada or USA"
  ]
}
\end{lstlisting}

All three suggestions reference this user's real \code{actual\_top} categories, not a paraphrase of
the stated query --- reproducing, with genuine model output, the same qualitative pattern already
observed with the mocked-LLM run.

\section{Interpretation and Limitations}
\label{sec:rec-limits}

The schema-grounding check (Section~\ref{sec:rec-validation}) is intentionally weak, and is reported
as such rather than presented as the precision filter originally planned. The candidate pool size
(five) is an arbitrary choice, not derived from data; a sensitivity check across pool sizes was not
performed. The behavioral-alignment ranking axis only functions for the Movie domain today, since
behavioral profiles only exist there (Section~\ref{sec:arch-vkgs}); it correctly returns $0.0$ when
no profile exists, so Tourism recommendations degrade to history-similarity-only ranking rather than
erroring, but this means the thesis's core stated-vs-actual contribution is only demonstrable
end-to-end on the Movie domain --- consistent with every earlier chapter's own scoping note. A
systematic, quantitative check of suggestion quality across many users and queries, rather than the
worked examples in Section~\ref{sec:rec-verification}, is the subject of Chapter~\ref{ch:eval}.
